% L12a jump hidden Markov model schematic: the N states ordered from the most
% negative to the most positive daily growth rate, each with its Student-t
% emission; ordinary days step with the counted transition matrix, and on a
% day with no episode under way a jump episode starts with probability epsilon,
% lasts a Poisson number of days with mean lambda, and draws each day's state
% from the most negative or the most positive states.
%
% Ported from the INFORMS 2026 talk (varnerlab/INFORMS-Annual-Meeting-2026-
% Varnerlab, talk/figs/jumphmm_architecture.tex, also Fig. jumphmm of the
% paper). Same layout; the labels use the lecture's notation (T_nu, hat P,
% epsilon, lambda, states 1..N) and the course palette: the most negative
% states in carnelian, the most positive in blue.
%
% The notebook SVG carries a dark palette added by theme_svg.py; keep new
% colors in that script's mapping when changing the palette here.
\documentclass[tikz,border=4pt]{standalone}
\usepackage{vnfigure}

\begin{document}
\begin{tikzpicture}[
  state/.style    = {circle, draw=vnink, line width=1.4pt, minimum size=1.8cm,
                     inner sep=0pt, font=\LARGE, text=vnink},
  tailbot/.style  = {state, draw=vncarnelian},
  tailtop/.style  = {state, draw=vnlink},
  midstate/.style = {state, draw=vnmuted},
  dots/.style     = {font=\Huge, text=vnmuted, yshift=-2pt},
  markov/.style   = {<->, >={Stealth[length=7pt, width=6pt]}, line width=0.9pt, vnmuted,
                     shorten <=3pt, shorten >=3pt},
  trigger/.style  = {->, >=Stealth, thick, dashed, vnink},
  jumparc/.style  = {->, >=Stealth, very thick},
  note/.style     = {font=\LARGE, fill=white, inner sep=2pt},
]
  % the chain: the most negative states, the k-th state, and the most positive states
  \node[tailbot]  (s1)   at ( 0.0, 0) {$1$};
  \node[tailbot]  (s2)   at ( 4.6, 0) {$2$};
  \node[dots]     (ld)   at ( 9.2, 0) {$\cdots$};
  \node[midstate] (sk)   at (13.8, 0) {$k$};
  \node[dots]     (rd)   at (18.4, 0) {$\cdots$};
  \node[tailtop]  (sNm1) at (23.0, 0) {$N\!-\!1$};
  \node[tailtop]  (sN)   at (27.6, 0) {$N$};
  \draw[markov] (s1)  -- (s2);
  \draw[markov] (s2)  -- (ld.west |- s2);
  \draw[markov] (ld.east |- sk) -- (sk);
  \draw[markov] (sk)  -- (rd.west |- sk);
  \draw[markov] (rd.east |- sNm1) -- (sNm1);
  \draw[markov] (sNm1)-- (sN);
  % each state's emission, a bell above it, on a thin base line
  \foreach \x/\c in {0.0/vncarnelian, 4.6/vncarnelian, 13.8/vnmuted, 23.0/vnlink, 27.6/vnlink} {
    \draw[\c, thick]
      plot[domain=-0.55:0.55, samples=40, smooth, variable=\t]
      ({\x+\t}, {1.3 + 0.85*exp(-\t*\t/0.04)});
    \draw[vnrule, thin] ({\x-0.55},1.3) -- ({\x+0.55},1.3);
  }
  % an episode starts from any state on a day with no episode under way; the dashed
  % trigger starts above the k-th emission so it does not cut the bell
  \node[draw=vnink, rounded corners=5pt, very thick,
        minimum width=2.8cm, minimum height=1.1cm,
        align=center, font=\LARGE, text=vnink] (clock) at (13.8, 3.9)
        {episode starts, prob.~$\epsilon$\\ lasts $\mathrm{Poisson}(\lambda)$ days};
  \draw[trigger] (13.8, 2.25) -- (clock.south)
    node[note, pos=0.5, right=5pt, anchor=west, text=vnmuted] {no episode under way};
  % each episode day draws its state from the most negative states (probability slightly
  % above one half) or the most positive states, chosen afresh every day
  \draw[jumparc, vncarnelian, rounded corners=10pt]
    (clock.west) --
    node[note, above, text=vncarnelian, pos=0.5, yshift=1pt]
      {each episode day, prob.~slightly above $\tfrac{1}{2}$}
    (-1.6, 3.9) -- (-1.6, 0) -- (s1.west);
  \draw[jumparc, vnlink, rounded corners=10pt]
    (clock.east) --
    node[note, above, text=vnlink, pos=0.5, yshift=1pt]
      {each episode day, prob.~slightly below $\tfrac{1}{2}$}
    (29.2, 3.9) -- (29.2, 0) -- (sN.east);
  % the tails, braced under their states, and the state index between them
  \draw[decorate, decoration={brace, amplitude=6pt, mirror}, vncarnelian, thick]
    (-0.8, -1.0) -- (5.4, -1.0)
    node[midway, below=5pt, text=vncarnelian, font=\LARGE] {most negative states};
  \draw[decorate, decoration={brace, amplitude=6pt, mirror}, vnlink, thick]
    (22.2, -1.0) -- (28.4, -1.0)
    node[midway, below=5pt, text=vnlink, font=\LARGE] {most positive states};
  \draw[->, >={Stealth[length=7pt, width=6pt]}, vnmuted, line width=0.9pt]
    (8.0, -1.5) -- (19.6, -1.5)
    node[midway, fill=white, inner sep=4pt, font=\LARGE, text=vnmuted] {higher growth rate};
  % what the k-th state emits, and the ordinary step
  \node[note, align=right, anchor=east, text=vnmuted]
    at (13.1, 1.75) {emission of state $k$\\[-1pt] $\mu_k + \sigma_k\,T_\nu$};
  \node[note, text=vnmuted, anchor=south west]
    at (14.75, 0.25) {step with $\hat{\mathbf{P}}$, prob.~$1-\epsilon$};
  % the box around what is drawn
  \pgfresetboundingbox
  \path[use as bounding box] (-1.75,-2.1) rectangle (29.35,4.95);
\end{tikzpicture}
\end{document}
